\chapter{Polinômios}\label{ch:b1:poly}

Os \omterm{def:b1:poly:def}{polinômios} são as funções prediletas do algebrista --- salvo que
aqui eles não são tratados como funções, e sim como expressões formais
numa indeterminada $X$, somadas e multiplicadas pelas regras de um
\omterm{def:b1:structures:ring}{anel} comutativo. A teoria corre de modo notavelmente paralelo ao
\cref{ch:b1:arith}: uma divisão euclidiana, um mdc e relações de Bézout,
elementos irredutíveis e uma fatoração única. Ao longo do capítulo, $K$
denota $\Q$, $\R$ ou $\C$.

\section{O anel $K[X]$}

\begin{definition}[Polinômio, grau]\label{def:b1:poly:def}
Um \emph{polinômio}\index{polinômio} com coeficientes em $K$ é uma
soma formal
\[
P = a_0 + a_1 X + a_2 X^2 + \dots + a_n X^n
= \sum_{k} a_k X^k,
\]
com $a_k \in K$ todos nulos a partir de um certo índice. Com a adição
natural e o produto
\[
\Bigl(\sum_i a_i X^i\Bigr)\Bigl(\sum_j b_j X^j\Bigr)
= \sum_k \Bigl(\sum_{i+j=k} a_i b_j\Bigr) X^k,
\]
o \omterm{def:b1:logic:sets}{conjunto} $K[X]$ é um \omterm{def:b1:structures:ring}{anel} comutativo. O \emph{grau} $\deg P$ de
$P \neq 0$ é o maior $n$ com $a_n \neq 0$; $a_n$ é o
\emph{coeficiente líder} ($P$ é \emph{mônico}\index{polinômio
mônico} quando $a_n = 1$) e, por convenção, $\deg 0 = -\infty$.
Todo polinômio define uma função $x \mapsto P(x)$ em $K$ por
substituição.
\end{definition}

\begin{proposition}[Regras dos graus; domínio de integridade]\label{prop:b1:poly:degree}
Para $P, Q \in K[X]$:
\[
\deg(P + Q) \leq \max(\deg P, \deg Q),
\qquad
\deg(PQ) = \deg P + \deg Q .
\]
Consequentemente, $K[X]$ é um \omterm{def:b1:structures:field}{domínio de integridade}, e as suas unidades
são as constantes não nulas.
\end{proposition}

\begin{proof}
A regra da soma é clara (os coeficientes acima do máximo se anulam).
Para o produto, sejam $a_m$ e $b_n$ os coeficientes líderes: o
coeficiente de $X^{m+n}$ em $PQ$ é $a_m b_n \neq 0$ ($K$ é um
\omterm{def:b1:structures:field}{corpo}, logo um \omterm{def:b1:structures:field}{domínio de integridade}), e todos os coeficientes
superiores se anulam. Se $P, Q \neq 0$, então $\deg PQ = \deg P + \deg Q \geq 0$, de modo que $PQ \neq
0$: \omterm{def:b1:structures:field}{domínio de integridade}. Se $PQ = 1$,
então $\deg P + \deg Q = 0$ força $\deg P = \deg Q = 0$: os elementos
invertíveis são as constantes invertíveis, isto é, todo o $K^*$.
\end{proof}

\begin{theorem}[Divisão euclidiana]\label{thm:b1:poly:division}
Sejam $A, B \in K[X]$ com $B \neq 0$. Existe exatamente um par $(Q,
R)$ de \omterm{def:b1:poly:def}{polinômios} com
\[
A = BQ + R, \qquad \deg R < \deg B .
\]
\end{theorem}

\begin{proof}
\emph{Existência}, por indução forte em $\deg A$. Se $\deg A < \deg
B$, tome $(Q, R) = (0, A)$. Caso contrário, escreva $A = a X^m + \dots$, $B =
b X^n + \dots$ com $m \geq n$; o \omterm{def:b1:poly:def}{polinômio} $A_1 = A - \frac ab
X^{m-n} B$ tem grau $< m$ (os termos líderes se cancelam), de modo que, por
indução, $A_1 = BQ_1 + R$ com $\deg R < \deg B$, e $A = B(Q_1 +
\frac ab X^{m-n}) + R$.

\emph{Unicidade}: se $BQ + R = BQ' + R'$, então $B(Q - Q') = R' - R$
com $\deg(R' - R) < \deg B$; pela regra dos graus, isso força $Q - Q'
= 0$ e, em seguida, $R = R'$.
\end{proof}

\begin{example}\label{ex:b1:poly:division}
Divida $A = X^4 + X^3 - 2X + 1$ por $B = X^2 + 1$:
\[
X^4 + X^3 - 2X + 1 = (X^2 + 1)(X^2 + X - 1) + (-3X + 2).
\]
(Cálculo: subtraia $X^2 B$, depois $X B$, depois $-B$; o resto
$-3X + 2$ tem grau $1 < 2$.)
\end{example}

\begin{method}[Esquema de Horner]\label{met:b1:poly:horner}
Para avaliar $P = a_nX^n + \dots + a_0$ em $x$, ou para dividir $P$ por
$X - x$, evite calcular potências: leia os coeficientes da esquerda para
a direita e itere \emph{multiplique por $x$, some o próximo
coeficiente}:
\[
b_n = a_n, \qquad b_{k} = a_{k} + x\,b_{k+1}
\quad (k = n-1, \dots, 0) .
\]
Então $b_0 = P(x)$ e os $b_k$ anteriores são os coeficientes do
quociente: $P = (X - x)(b_nX^{n-1} + \dots + b_1) + b_0$
(expanda e compare). Exemplo: $P = X^4 - 5X^3 + 6X^2 + 4X - 8$
em $x = 2$: os $b$ valem $1, -3, 0, 4, 0$, de modo que $P(2) = 0$ e $P
= (X-2)(X^3 - 3X^2 + 4)$ --- uma linha em vez de uma longa
divisão, e $n$ multiplicações em vez das $\approx n^2/2$
da avaliação ingênua. Iterar o esquema no mesmo ponto
extrai as \omterm{def:b1:poly:derivative}{multiplicidades} (compare com o
\cref{ex:b1:poly:multexample}).
\end{method}

\begin{remark}[{A aritmética de $K[X]$}]
Com a divisão euclidiana em mãos, toda a aritmética do
\cref{ch:b1:arith} se transfere para $K[X]$, com as mesmas demonstrações,
sendo o grau o papel do valor absoluto: mdc (normalizado para ser
\omterm{def:b1:poly:def}{mônico}), \omterm{met:b1:arith:euclid}{algoritmo de Euclides} estendido, identidade de Bézout, lema de
Gauss, \omterm{def:b1:poly:def}{polinômios} irredutíveis e fatoração única. Usamos livremente esses
resultados transferidos, e o \cref{exo:b1:poly:6} reencena um deles.
\end{remark}

\section{Raízes}

\begin{theorem}[Teorema do fator]\label{thm:b1:poly:factor}
Sejam $P \in K[X]$ e $a \in K$. O resto de $P$ na divisão por
$X - a$ é a constante $P(a)$. Em particular,
\[
P(a) = 0 \iff (X - a) \mid P .
\]
Mais geralmente, raízes distintas $a_1, \dots, a_r$ de $P$ dão a
fatoração $P = (X - a_1)\cdots(X - a_r)\, Q$.
\end{theorem}

\begin{proof}
Divida: $P = (X - a) Q + R$ com $\deg R < 1$, de modo que $R$ é uma constante
$c$; substituindo $X = a$ (a substituição respeita somas e produtos),
obtém-se $P(a) = c$. A equivalência segue. Para várias raízes,
faça indução em $r$: o caso $r = 1$ é a equivalência recém-demonstrada.
Suponha o enunciado para $r - 1$ raízes e sejam $a_1, \dots, a_r$
raízes distintas de $P$. Escreva $P = (X - a_1)Q_1$; para cada $i
\geq 2$, substituindo $a_i$:
\[
0 = P(a_i) = (a_i - a_1)\,Q_1(a_i),
\qquad a_i - a_1 \neq 0 ,
\]
e, como $K$ não tem divisores de zero, $Q_1(a_i) = 0$: os $r - 1$
pontos distintos $a_2, \dots, a_r$ são raízes de $Q_1$. A
hipótese de indução fatora $Q_1 = (X - a_2)\cdots(X - a_r)\,Q$,
e substituir de volta dá a afirmação.
\end{proof}

\begin{corollary}[Um polinômio de grau $n$ tem no máximo $n$ raízes]\label{cor:b1:poly:nroots}
Um $P \in K[X]$ não nulo de grau $n$ tem no máximo $n$ raízes distintas
em $K$. Consequentemente, um \omterm{def:b1:poly:def}{polinômio} (de grau $\leq n$) que se anula em
$n + 1$ pontos distintos é o \omterm{def:b1:poly:def}{polinômio} nulo, e dois \omterm{def:b1:poly:def}{polinômios}
de grau $\leq n$ que coincidem em $n+1$ pontos são iguais.
\end{corollary}

\begin{proof}
Se $a_1, \dots, a_r$ são raízes distintas, o \cref{thm:b1:poly:factor}
dá $P = (X-a_1)\cdots(X-a_r) Q$, de modo que $n = \deg P \geq r$. As duas
consequências decorrem por absurdo e por diferença.
\end{proof}

\begin{example}[O truque do polinômio auxiliar]\label{ex:b1:poly:auxiliary}
Seja $P$ o \omterm{def:b1:poly:def}{polinômio} de grau $\leq n$ com
\[
P(k) = \frac{k}{k+1} \qquad (k = 0, 1, \dots, n) ;
\]
ele existe e é único pela \omterm{thm:b1:poly:lagrange}{interpolação de Lagrange}, abaixo. Quanto vale
$P(n+1)$? Elimine os denominadores: o \omterm{def:b1:poly:def}{polinômio} $Q = (X+1)P - X$ tem
grau $\leq n + 1$ e se anula nos $n + 1$ pontos $0, 1,
\dots, n$, de modo que, pelo \cref{thm:b1:poly:factor},
\[
Q = c\,X(X-1)(X-2)\cdots(X-n)
\]
para alguma constante $c$. Avalie onde $Q$ é conhecido de modo
independente: em $X = -1$, $Q(-1) = 0 \cdot P(-1) + 1 = 1$, ao passo que o produto
vale $(-1)(-2)\cdots(-1-n) = (-1)^{n+1}(n+1)!$; portanto, $c =
\frac{(-1)^{n+1}}{(n+1)!}$. Agora avalie em $X = n + 1$:
\[
(n+2)\,P(n+1) - (n+1) = Q(n+1) = c\,(n+1)! = (-1)^{n+1} ,
\]
de modo que $P(n+1) = \dfrac{(n+1) + (-1)^{n+1}}{n+2}$: igual a $1$ para
$n$ ímpar, e a $\frac{n}{n+2}$ para $n$ par --- o
\omterm{def:b1:poly:def}{polinômio} interpolador \emph{não} continua o padrão
$\frac{n+1}{n+2}$. O truque a reter: codifique os dados como
raízes de um \omterm{def:b1:poly:def}{polinômio} auxiliar, identifique a constante desconhecida
num ponto fora dos dados e colha o resultado.
\end{example}

\begin{definition}[Derivada, multiplicidade]\label{def:b1:poly:derivative}
A \emph{derivada formal} de $P = \sum a_k X^k$ é $P' = \sum_{k
\geq 1} k\,a_k X^{k-1}$; ela satisfaz as regras usuais $(P+Q)' = P' +
Q'$, $(PQ)' = P'Q + PQ'$ (verificadas nos monômios e estendidas por
linearidade). Uma raiz $a$ de $P$ tem
\emph{multiplicidade}\index{multiplicidade de uma raiz} $m \geq 1$ quando
$(X-a)^m \mid P$, mas $(X-a)^{m+1} \nmid P$; a raiz é \emph{simples}
se $m = 1$, e \emph{múltipla} se $m \geq 2$.
\end{definition}

\begin{proposition}[Multiplicidade via derivadas]\label{prop:b1:poly:multiplicity}
$a$ é raiz de $P$ de \omterm{def:b1:poly:derivative}{multiplicidade} $\geq m$ se, e somente se,
\[
P(a) = P'(a) = \dots = P^{(m-1)}(a) = 0 .
\]
Em particular, $a$ é raiz múltipla de $P$ se, e somente se, $P(a) =
P'(a) = 0$.
\end{proposition}

\begin{proof}
Escreva $P = (X - a)^m Q + R$, em que $R$ é o resto da divisão
por $(X-a)^m$, $\deg R < m$. Derivando $k \leq m - 1$ vezes e
avaliando em $a$: o primeiro termo contribui com $0$ (cada derivada
retém um fator $(X-a)$), de modo que $P^{(k)}(a) = R^{(k)}(a)$.

Ora, um \omterm{def:b1:poly:def}{polinômio} $R$ de grau $< m$ fica determinado por $R(a), R'(a),
\dots, R^{(m-1)}(a)$: escrevendo $R = \sum_{k < m} c_k (X - a)^k$
(o que é possível: expanda as potências de $X = (X - a) + a$), encontra-se
$R^{(k)}(a) = k!\, c_k$. Portanto: todos os $P^{(k)}(a) = 0$ para $k < m$
$\iff$ todos os $c_k = 0$ $\iff$ $R = 0$ $\iff$ $(X-a)^m \mid P$.
\end{proof}

\begin{example}[Calculando uma multiplicidade]\label{ex:b1:poly:multexample}
Qual é a \omterm{def:b1:poly:derivative}{multiplicidade} da raiz $2$ em $P = X^4 - 5X^3 + 6X^2
+ 4X - 8$? Avalie as derivadas sucessivas em $2$:
\[
P(2) = 16 - 40 + 24 + 8 - 8 = 0, \qquad
P'(2) = 32 - 60 + 24 + 4 = 0,
\]
\[
P''(2) = 48 - 60 + 12 = 0, \qquad
P'''(2) = 48 - 30 = 18 \neq 0
\]
(com $P' = 4X^3 - 15X^2 + 12X + 4$, $P'' = 12X^2 - 30X + 12$,
$P''' = 24X - 30$). Três valores nulos e depois um não nulo:
\omterm{def:b1:poly:derivative}{multiplicidade} exatamente $3$. Dividindo, $P = (X - 2)^3(X + 1)$ ---
o que se confere expandindo $(X-2)^3 = X^3 - 6X^2 + 12X - 8$ e
multiplicando por $X + 1$. A ideia: as \omterm{def:b1:poly:derivative}{multiplicidades} são lidas em
\emph{avaliações}, sem necessidade de fatoração --- e é exatamente
assim que se as detecta quando a fatoração está fora de alcance.
\end{example}

\begin{example}[Detectando raízes múltiplas com um mdc]\label{ex:b1:poly:gcdmultiple}
Quando nenhuma raiz é conhecida, a \cref{prop:b1:poly:multiplicity} ainda
fornece um detector \emph{global} de raízes múltiplas: $a$ é raiz
múltipla de $P$ se, e somente se, é raiz comum de $P$ e $P'$, de modo
que $P$ tem raiz múltipla (em $\C$) se, e somente se, $\gcd(P, P') \neq 1$ ---
o que é computável pelo \omterm{met:b1:arith:euclid}{algoritmo de Euclides} sem resolver nada.
Exemplo: $P =
X^3 - 3X + 2$, $P' = 3X^2 - 3 = 3(X - 1)(X + 1)$. Testando as
raízes $\pm1$ de $P'$ dentro de $P$: $P(1) = 0$, mas $P(-1) = 4$, logo
\[
\gcd(P, P') = X - 1 :
\]
a raiz $1$ é múltipla; dividindo duas vezes, $P = (X - 1)^2(X + 2)$.
O mdc chega a relatar o \omterm{def:b1:logic:sets}{conjunto} completo das raízes múltiplas, cada uma
com \omterm{def:b1:poly:derivative}{multiplicidade} abaixada de uma unidade --- fato que todo sistema de
álgebra computacional explora para ``fatorar sem quadrados'' antes de
qualquer caça a raízes, e o gêmeo polinomial dos argumentos de ausência
de raiz múltipla do \cref{exo:b1:poly:9}.
\end{example}

\begin{theorem}[Teorema fundamental da álgebra]\label{thm:b1:poly:fta}
Todo \omterm{def:b1:poly:def}{polinômio} não constante de $\C[X]$ tem uma raiz em
$\C$.\index{teorema fundamental da álgebra}
\end{theorem}
\admitted

\begin{remark}
Apesar do nome, o teorema é um enunciado de \emph{análise}:
toda demonstração conhecida usa a completude de $\R$ de alguma forma, e
nenhuma é puramente algébrica --- a demonstração honesta é dada no
volume do terceiro ano de graduação, uma vez disponíveis a integração
complexa ou argumentos de compacidade. O que este capítulo genuinamente
demonstra é a \emph{redução}: concedida uma raiz para todo \omterm{def:b1:poly:def}{polinômio}
não constante, as fatorações completas sobre $\C$ e sobre $\R$ abaixo
decorrem por pura álgebra.
\end{remark}

\begin{corollary}[Fatoração sobre $\C$ e sobre $\R$]\label{cor:b1:poly:factorization}
\begin{enumerate}
  \item Todo $P \in \C[X]$ não nulo se fatora como
        \[
        P = c\, (X - a_1)^{m_1} \cdots (X - a_r)^{m_r},
        \]
        com $c$ o coeficiente líder, $a_i$ as raízes complexas
        distintas e $\sum m_i = \deg P$: \emph{contadas com \omterm{def:b1:poly:derivative}{multiplicidade},
        um \omterm{def:b1:poly:def}{polinômio} de grau $n$ tem exatamente $n$ raízes complexas}.
  \item Todo $P \in \R[X]$ não nulo se fatora sobre $\R$ como
        \[
        P = c \prod_i (X - a_i)^{m_i} \prod_j (X^2 + p_j X +
        q_j)^{n_j},
        \]
        sendo os fatores quadráticos distintos e com $p_j^2 - 4q_j < 0$
        (sem raízes reais).
\end{enumerate}
\end{corollary}

\begin{proof}
(1) Indução no grau, destacando uma raiz de cada vez pelo
\cref{thm:b1:poly:factor}; a contagem dos graus fecha em cada passo.

(2) Seja $P$ com coeficientes reais. Se $z$ é raiz complexa de
\omterm{def:b1:poly:derivative}{multiplicidade} $m$, então $\conj z$ também é: conjugar $P(z) = 0$ dá
$P(\conj z) = \conj{P(z)} = 0$ (os coeficientes são os seus próprios
conjugados), e o mesmo se aplica às derivadas
(\cref{prop:b1:poly:multiplicity}). Agrupe as raízes não reais em
pares conjugados: cada par contribui com
\[
(X - z)(X - \conj z) = X^2 - 2\Re(z)\, X + \abs z^2 ,
\]
uma quadrática real de discriminante negativo. As raízes reais contribuem
com os fatores lineares.
\end{proof}

\begin{example}\label{ex:b1:poly:factorexample}
$X^4 + 4$ foi fatorado sobre $\R$ no \cref{exo:b1:complex:5}, emparelhando
as quatro raízes complexas $\pm 1 \pm \iu$:
$X^4 + 4 = (X^2 - 2X + 2)(X^2 + 2X + 2)$. Nenhuma das quadráticas se decompõe
sobre $\R$ (discriminantes $-4$). Note: um \omterm{def:b1:poly:def}{polinômio} real
\emph{irredutível} tem grau $1$ ou $2$ --- é exatamente o que o
teorema de fatoração diz. O mesmo emparelhamento de conjugados aplicado a
$X^4 + 1$, cujas raízes são $\eu^{\pm\iu\pi/4}$ e
$\eu^{\pm3\iu\pi/4}$: cada par contribui com $X^2 -
2\cos\theta\,X + 1$, de modo que
\[
X^4 + 1 = \bigl(X^2 - \sqrt2\,X + 1\bigr)
\bigl(X^2 + \sqrt2\,X + 1\bigr) ,
\]
uma identidade invisível a tentativas ingênuas de fatoração sobre $\Q$
--- o preço de insistir em coeficientes reais (aqui, até irracionais)
e um insumo padrão para integrar $\frac1{x^4 +
1}$ no \cref{ch:b1:integration}.
\end{example}

\begin{omfigure}
\begin{tikzpicture}[xscale=2.6, yscale=1.35]
  \draw[->, thin] (-1.15,0) -- (1.2,0) node[below] {$x$};
  \draw[->, thin] (0,-1.35) -- (0,1.4) node[left] {$y$};
  \draw[dashed, gray] (-1,1) -- (1,1);
  \draw[dashed, gray] (-1,-1) -- (1,-1);
  \draw[gray, thin] (-1,-1.1) -- (-1,1.1) (1,-1.1) -- (1,1.1);
  \draw[omThm, very thick, domain=-1:1, samples=300, smooth]
    plot (\x, {16*\x^5 - 20*\x^3 + 5*\x});
  \node[right, font=\small] at (1.02,1) {$1$};
  \node[right, font=\small] at (1.02,-1) {$-1$};
  \fill[omThm] (-1,-1) circle (0.018) (-0.809,1) circle (0.018)
    (-0.309,-1) circle (0.018) (0.309,1) circle (0.018)
    (0.809,-1) circle (0.018) (1,1) circle (0.018);
\end{tikzpicture}

{\small O \omterm{def:b1:poly:def}{polinômio} de Chebyshev $T_5 = 16X^5 - 20X^3 + 5X$ em
$\intcc{-1}1$: ele oscila exatamente entre $-1$ e $1$,
tocando os limites em seis pontos (marcados). Essa
\emph{equioscilação} é o que faz de $2^{-4}T_5$ a quíntica
\omterm{def:b1:poly:def}{mônica} de menor norma do supremo no intervalo
(\cref{exo:b1:poly:10} e problema de fim de semana).}
\end{omfigure}

\begin{remark}[Armadilhas frequentes com polinômios]\label{rem:b1:poly:pitfalls}
\begin{enumerate}
  \item \emph{\omterm{def:b1:poly:def}{Polinômio} versus função.} Sobre $K = \Q, \R, \C$
        as duas noções coincidem (funções iguais têm coeficientes
        iguais, pelo \cref{cor:b1:poly:nroots} e pela
        infinitude de $K$), mas conceitualmente um \omterm{def:b1:poly:def}{polinômio} é
        a sua lista de coeficientes: sobre o \omterm{def:b1:structures:field}{corpo} de dois elementos
        $\Z/2\Z$ do \cref{ch:b1:structures}, $X^2 + X$ se anula
        nos dois pontos e, no entanto, não é o \omterm{def:b1:poly:def}{polinômio} nulo.
  \item \emph{Graus sob adição.} $\deg(P + Q)$ pode cair
        abaixo de $\max(\deg P, \deg Q)$ quando os termos líderes se cancelam;
        escrever ``$\deg(P + Q) = \max(\dots)$'' só é seguro para
        graus distintos.
  \item \emph{Contar as raízes corretamente.} ``$n$ raízes'' no
        \cref{cor:b1:poly:factorization} significa \emph{com
        \omterm{def:b1:poly:derivative}{multiplicidade}, em $\C$}: $X^2 + 1$ não tem raízes reais, e
        $(X-1)^2$ tem uma raiz distinta, mas duas com
        \omterm{def:b1:poly:derivative}{multiplicidade}. Enunciados que misturam as três contagens são a
        fonte mais comum de demonstrações falsas.
  \item \emph{A irredutibilidade depende do \omterm{def:b1:structures:field}{corpo}.} $X^2 - 2$ é
        irredutível sobre $\Q$ e se decompõe sobre $\R$; $X^2 + 1$ é
        irredutível sobre $\R$ e se decompõe sobre $\C$. A palavra
        solta ``irredutível'' nada significa enquanto não se nomear o
        \omterm{def:b1:structures:field}{corpo} dos coeficientes.
\end{enumerate}
\end{remark}

\section{Coeficientes e raízes}

\begin{theorem}[Relações de Girard]\label{thm:b1:poly:vieta}
Seja $P = X^n + c_{n-1} X^{n-1} + \dots + c_0$ \omterm{def:b1:poly:def}{mônico}, com raízes
$a_1, \dots, a_n \in \C$ (com \omterm{def:b1:poly:derivative}{multiplicidade}). Então\index{relações de
Girard}
\[
\sum_i a_i = -c_{n-1},
\qquad
\sum_{i < j} a_i a_j = c_{n-2},
\qquad \dots, \qquad
a_1 a_2 \cdots a_n = (-1)^n c_0 ,
\]
sendo a $k$-ésima função simétrica das raízes igual a $(-1)^k c_{n-k}$.
\end{theorem}

\begin{proof}
Pelo \cref{cor:b1:poly:factorization}, $P = (X - a_1)\cdots(X -
a_n)$ (\omterm{def:b1:poly:def}{mônico}, todas as raízes listadas). Expandir o produto
distributivamente produz um termo por escolha, em cada
fator, de $X$ ou do termo de raiz $-a_i$: escolher as raízes nos
fatores de índices $i_1 < \dots < i_k$ e $X$ nos $n - k$
restantes contribui com $(-a_{i_1})\cdots(-a_{i_k})\,X^{n-k}$. Agrupando
pela potência de $X$:
\[
P = \sum_{k=0}^{n} (-1)^k
\Bigl(\sum_{i_1 < \dots < i_k} a_{i_1}\cdots a_{i_k}\Bigr)
X^{n-k} ,
\]
e identificar com $P = \sum_k c_{n-k}X^{n-k}$ (os coeficientes
são únicos, \cref{def:b1:poly:def}) dá $c_{n-k} = (-1)^k
\sigma_k$, isto é, $\sigma_k = (-1)^kc_{n-k}$, em que $\sigma_k$
denota a $k$-ésima função simétrica exibida acima. Os três
casos exibidos são $k = 1$, $k = 2$ e $k = n$.
\end{proof}

\begin{example}\label{ex:b1:poly:vieta}
Para a quadrática $X^2 - sX + p$: soma das raízes $s$, produto $p$ ---
já usados repetidamente (\cref{exo:b1:complex:8}). Para uma cúbica \omterm{def:b1:poly:def}{mônica}
$X^3 + aX^2 + bX + c$ com raízes $\alpha, \beta, \gamma$:
\[
\alpha + \beta + \gamma = -a,
\quad
\alpha\beta + \beta\gamma + \gamma\alpha = b,
\quad
\alpha\beta\gamma = -c ,
\]
o que permite calcular grandezas simétricas como $\alpha^2 + \beta^2
+ \gamma^2 = a^2 - 2b$ sem resolver nada.
\end{example}

\begin{example}[Transformar as raízes sem encontrá-las]\label{ex:b1:poly:transformroots}
Sejam $\alpha, \beta$ as raízes de $X^2 - 3X + 1$. Que quadrática \omterm{def:b1:poly:def}{mônica}
tem raízes $\alpha^2, \beta^2$? Por Girard, $\alpha +
\beta = 3$ e $\alpha\beta = 1$, de modo que
\[
\alpha^2 + \beta^2 = (\alpha+\beta)^2 - 2\alpha\beta = 7,
\qquad
\alpha^2\beta^2 = (\alpha\beta)^2 = 1 :
\]
a resposta é $X^2 - 7X + 1$ --- obtida sem calcular
$\alpha = \frac{3 + \sqrt5}2$. (Verificação: $\alpha^2 = \frac{7 +
3\sqrt5}2$, e de fato $\alpha^2 + \beta^2 = 7$.) A mesma
estratégia trata os inversos (transformações do tipo $X^2 - \frac ba X + \frac
ca$), translações e quaisquer dados simétricos: as \omterm{thm:b1:poly:vieta}{relações de
Girard} convertem perguntas sobre raízes \emph{desconhecidas} em álgebra
sobre coeficientes \emph{conhecidos}. Elas servirão constantemente
quando as raízes forem autovalores (\cref{ch:b1:det}).
\end{example}

\begin{example}[Equações recíprocas]\label{ex:b1:poly:palindrome}
Resolva $X^4 + X^3 - 4X^2 + X + 1 = 0$. Os coeficientes se leem do
mesmo modo nos dois sentidos, de sorte que $0$ não é raiz e dividir por
$X^2$ não perde soluções:
\[
X^2 + X - 4 + \frac1X + \frac1{X^2} = 0 .
\]
Ponha $y = X + \frac1X$: então $X^2 + \frac1{X^2} = y^2 - 2$, e
a equação colapsa em
\[
y^2 + y - 6 = 0 \iff (y + 3)(y - 2) = 0 .
\]
Desdobre cada valor por meio de $X^2 - yX + 1 = 0$: para $y = 2$, $X^2 -
2X + 1 = (X - 1)^2$ dá a raiz dupla $1$; para $y = -3$,
$X^2 + 3X + 1 = 0$ dá $X = \frac{-3 \pm \sqrt5}2$. Quatro raízes
com \omterm{def:b1:poly:derivative}{multiplicidade} para uma quártica, como o
\cref{cor:b1:poly:factorization} exige --- obtidas resolvendo
duas quadráticas. O truque cobre todo \omterm{def:b1:poly:def}{polinômio} \emph{recíproco}:
as suas raízes vêm em pares inversos $\{x, 1/x\}$
(substitua $X$ por $1/X$ e elimine denominadores), e $y = X +
\frac1X$ é precisamente a grandeza constante em tais pares,
reduzindo o grau à metade.
\end{example}

\begin{theorem}[Interpolação de Lagrange]\label{thm:b1:poly:lagrange}
Sejam $x_0, \dots, x_n$ pontos distintos de $K$ e $y_0, \dots, y_n
\in K$. Existe exatamente um $P \in K[X]$ de grau $\leq n$ com
$P(x_i) = y_i$ para todo $i$, a saber\index{interpolação de Lagrange}
\[
P = \sum_{i=0}^{n} y_i\, L_i,
\qquad
L_i = \prod_{j \neq i} \frac{X - x_j}{x_i - x_j} .
\]
\end{theorem}

\begin{proof}
Cada $L_i$ tem grau $n$ e satisfaz $L_i(x_i) = 1$, $L_i(x_j) =
0$ para $j \neq i$ (cada fator se anula no $x_j$ correspondente).
Logo, o $P$ exibido tem grau $\leq n$ e interpola. Unicidade:
dois \omterm{def:b1:poly:def}{polinômios} interpoladores de grau $\leq n$ coincidem nos $n+1$
pontos $x_i$ e, portanto, são iguais (\cref{cor:b1:poly:nroots}).
\end{proof}

\begin{remark}[Interlúdio: polinômios também são vetores]\label{rem:b1:poly:interlude}
Uma mudança de ponto de vista que o \cref{ch:b1:vspaces} tornará
oficial: os \omterm{def:b1:poly:def}{polinômios} de grau $\leq n$ formam um espaço em
que a adição e a multiplicação por escalar se comportam exatamente como
coordenadas --- um \omterm{def:b1:poly:def}{polinômio} \emph{é} a sua lista de $n + 1$
coeficientes. Três enunciados deste capítulo são álgebra linear
disfarçada. A \omterm{thm:b1:poly:lagrange}{interpolação de Lagrange} (\cref{thm:b1:poly:lagrange}) diz que
os dados de avaliação $(P(x_0), \dots, P(x_n))$ determinam $P$
de modo único: avaliar em $n + 1$ pontos é uma bijeção linear,
e os $L_i$ são a base adaptada a ela. A expansão $R =
\sum c_k(X - a)^k$ na demonstração da
\cref{prop:b1:poly:multiplicity} diz que as potências de
$(X - a)$ formam outro sistema de coordenadas, com $c_k =
R^{(k)}(a)/k!$ como coordenadas. E o
\cref{cor:b1:poly:nroots} --- mais raízes que o grau força
o \omterm{def:b1:poly:def}{polinômio} nulo --- é o motor de todas as unicidades: ele se tornará
``uma \omterm{def:b1:logic:map}{aplicação} linear \omterm{def:b1:logic:inj}{injetiva} num espaço de dimensão $n + 1$''
no \cref{ch:b1:findim}. Quando esses capítulos chegarem, o espaço
$K_n[X]$ será o seu exemplo predileto; vale chegar lá já
fluente nele.
\end{remark}

\begin{remark}[Onde este capítulo é usado]\label{rem:b1:poly:whereused}
A fatoração sobre $\R$ e $\C$
(\cref{cor:b1:poly:factorization}) é o motor das frações
parciais no \cref{ch:b1:fractions}, logo de uma vasta classe de
integrais no \cref{ch:b1:integration}. A expansão de um
\omterm{def:b1:poly:def}{polinômio} em potências de $(X - a)$, encontrada na demonstração da
\cref{prop:b1:poly:multiplicity}, é a sombra algébrica das
fórmulas de Taylor do \cref{ch:b1:taylor}. Os \omterm{def:b1:poly:def}{polinômios} característicos
já apareceram para as equações diferenciais (\cref{ch:b1:diffeq})
e voltam para as matrizes no \cref{ch:b1:det}; a \omterm{thm:b1:poly:lagrange}{interpolação
de Lagrange} é o primeiro teorema de existência e unicidade da
análise numérica, e os \omterm{pb:b1:poly:1}{polinômios de Chebyshev} do
\cref{exo:b1:poly:10} --- cuja otimalidade o problema de fim de semana
abaixo estabelece --- dizem a essa disciplina \emph{onde}
interpolar. Por fim, toda a aritmética de $K[X]$, copiada do
\cref{ch:b1:arith}, alimenta o estudo dos ideais de $K[X]$ e dos anéis
quocientes no volume do segundo ano de graduação.
\end{remark}

\section{Exercícios}

\begin{exercise}[$\star$]\label{exo:b1:poly:1}
Efetue as divisões euclidianas: $X^5 - 1$ por $X^2 + X + 1$; depois
$2X^4 + X^3 - X + 3$ por $X^2 - 2$.
\end{exercise}

\begin{exercise}[$\star$]\label{exo:b1:poly:2}
Para quais $n \in \N$ o \omterm{def:b1:poly:def}{polinômio} $X^2 + X + 1$ \omterm{def:b1:arith:divides}{divide} $X^{2n} + X^n + 1$?
\emph{Sugestão: as raízes de $X^2 + X + 1$ são $j$ e $j^2$, com $j =
\eu^{2\iu\pi/3}$; discuta $n$ módulo $3$.}
\end{exercise}

\begin{exercise}[$\star$]\label{exo:b1:poly:3}
Determine os reais $a, b$ para que $(X-1)^2$ divida $P = X^4 + aX^3 +
bX^2 + 1$, e depois fatore $P$ sobre $\R$ para esses valores.
\end{exercise}

\begin{exercise}[$\star$]\label{exo:b1:poly:4}
Fatore sobre $\C$ e sobre $\R$: $X^3 - 1$; $\;X^4 + X^2 + 1$;
$\;X^6 - 1$.
\end{exercise}

\begin{exercise}[$\star\star$]\label{exo:b1:poly:5}
Seja $P = X^3 - 6X^2 + 11X - 6$.
\begin{enumerate}
  \item Encontre as raízes racionais \emph{(uma raiz racional $p/q$ na
        forma irredutível de um \omterm{def:b1:poly:def}{polinômio} inteiro \omterm{def:b1:poly:def}{mônico} é um inteiro
        que \omterm{def:b1:arith:divides}{divide} o termo constante --- demonstre-o)} e fatore $P$.
  \item Sem resolver, calcule a soma dos quadrados e a soma
        dos inversos das raízes pelas \omterm{thm:b1:poly:vieta}{relações de Girard}, e confira na
        fatoração.
\end{enumerate}
\end{exercise}

\begin{exercise}[$\star\star$]\label{exo:b1:poly:6}
Calcule $\gcd(X^4 - 1,\; X^3 - X^2 + X - 1)$ pelo \omterm{met:b1:arith:euclid}{algoritmo de Euclides} e
escreva-o como combinação $AU + BV$ dos dois \omterm{def:b1:poly:def}{polinômios}.
\end{exercise}

\begin{exercise}[$\star\star$]\label{exo:b1:poly:7}
Seja $P \in \R[X]$ com $P(x) \geq 0$ para todo $x \in \R$. Demonstre que
$P$ é uma soma de dois quadrados de \omterm{def:b1:poly:def}{polinômios} reais: $P = A^2 + B^2$.
\emph{Sugestão: na fatoração real, as raízes reais têm \omterm{def:b1:poly:derivative}{multiplicidade}
par; escreva os fatores quadráticos como $(X - z)(X - \conj z)$
e use $\abs{\,\cdot\,}^2 = (\Re)^2 + (\Im)^2$ no produto dos
$(X - z)$.}
\end{exercise}

\begin{exercise}[$\star\star$]\label{exo:b1:poly:8}
Encontre o \omterm{def:b1:poly:def}{polinômio} $P$ de grau $\leq 2$ com $P(0) = 1$, $P(1) =
3$, $P(2) = 2$, primeiro pela fórmula de Lagrange e depois resolvendo o
sistema linear nos coeficientes. Verifique que as duas respostas
coincidem.
\end{exercise}

\begin{exercise}[$\star\star$]\label{exo:b1:poly:9}
Demonstre que $P = X^{2n+1} - 1$ tem exatamente uma raiz real e que, para
todo $n \geq 1$, o \omterm{def:b1:poly:def}{polinômio} $1 + X + \frac{X^2}{2!} + \dots +
\frac{X^n}{n!}$ não tem raiz múltipla \emph{(compare $P$ e
$P'$)}.
\end{exercise}

\begin{exercise}[$\star\star\star$]\label{exo:b1:poly:10}
(\omterm{pb:b1:poly:1}{Polinômios de Chebyshev}) Defina $T_0 = 1$, $T_1 = X$ e $T_{n+1} =
2X\,T_n - T_{n-1}$.
\begin{enumerate}
  \item Demonstre por indução que $T_n(\cos\theta) = \cos n\theta$ para
        todo $\theta$.
  \item Deduza as $n$ raízes de $T_n$ e o seu coeficiente líder.
  \item Demonstre que $\sup_{x \in \intcc{-1}{1}} \abs{T_n(x)} = 1$,
        atingido em $n + 1$ pontos de $\intcc{-1}{1}$.
\end{enumerate}
\end{exercise}

\begin{exercise}[$\star\star\star$]\label{exo:b1:poly:11}
Seja $P \in \C[X]$ não constante, com raízes distintas $a_1, \dots,
a_r$ (de \omterm{def:b1:poly:derivative}{multiplicidades} $m_1, \dots, m_r$). Demonstre a identidade
de funções racionais
\[
\frac{P'(X)}{P(X)} = \sum_{i=1}^{r} \frac{m_i}{X - a_i},
\]
e deduza o teorema de Gauss--Lucas: toda raiz de $P'$ está na
envoltória convexa das raízes de $P$ \emph{(avalie a identidade numa
raiz $w$ de $P'$ que não seja raiz de $P$, tome conjugados e
leia o resultado como sendo $w$ uma média ponderada dos $a_i$)}.
\end{exercise}

\begin{exercise}[$\star\star$]\label{exo:b1:poly:12}
(Filtro das \omterm{def:b1:complex:unity}{raízes da unidade}) Sejam $n \in \N^*$ e $j = \eu^{2\iu\pi/3}$.
Avaliando $(1 + X)^n$ em $1$, $j$ e $j^2$, demonstre que
\[
\sum_{k \geq 0} \binom{n}{3k}
= \frac{2^n + 2\cos\frac{n\pi}{3}}{3} ,
\]
e confira a fórmula para $n = 3$ e $n = 6$. \emph{Sugestão: $1 +
j^m + j^{2m}$ vale $3$ se $3 \mid m$ e $0$ caso contrário; e $1 +
j = \eu^{\iu\pi/3}$.}
\end{exercise}

\section{Problema: Polinômios de Chebyshev e o polinômio mais achatado}

\begin{problem}\label{pb:b1:poly:1}
Entre todos os \omterm{def:b1:poly:def}{polinômios} \emph{\omterm{def:b1:poly:def}{mônicos}} de grau $n$, qual deles fica
mais próximo de zero em $\intcc{-1}1$? A resposta --- o teorema de
Chebyshev, a certidão de nascimento da teoria da aproximação --- é
$2^{1-n}T_n$, em que $T_n$ é o \omterm{def:b1:poly:def}{polinômio} de Chebyshev do
\cref{exo:b1:poly:10}, e nenhum concorrente \omterm{def:b1:poly:def}{mônico} consegue superar o seu
desvio $2^{1-n}$.\index{polinômios de Chebyshev} Este problema
desenvolve a álgebra da família $(T_n)$ (\omterm{def:b1:structures:law}{lei de composição},
coeficientes explícitos, a família de segunda espécie $U_n$, uma
equação diferencial), demonstra o teorema de extremalidade com o seu
caso de igualdade e recolhe aplicações: nós ótimos de interpolação,
o valor exato de $\cos 36^\circ$ e uma \omterm{def:b1:arith:congruence}{congruência}
$T_p \equiv X^p \pmod p$. Ao longo do problema, $T_0 = 1$, $T_1 = X$,
$T_{n+1} = 2X\,T_n - T_{n-1}$, e usamos livremente
$T_n(\cos\theta) = \cos n\theta$ do \cref{exo:b1:poly:10}.

\medskip
\noindent\textbf{Parte I --- A família $(T_n)$.}
\begin{enumerate}
  \item Calcule $T_2, T_3, T_4, T_5$ pela recorrência. (Compare
        $T_3$ com a identidade $\cos3\theta = 4\cos^3\theta -
        3\cos\theta$ do \cref{ex:b1:complex:cos3}.)
  \item Demonstre por indução que $\deg T_n = n$, com coeficiente
        líder $2^{n-1}$ para $n \geq 1$, e que $T_n$ tem a
        paridade de $n$ (só aparecem potências pares ou só ímpares).
  \item Demonstre o princípio de unicidade: $T_n$ é o \emph{único}
        \omterm{def:b1:poly:def}{polinômio} que satisfaz $P(\cos\theta) = \cos n\theta$ para
        todo $\theta$. (Dois \omterm{def:b1:poly:def}{polinômios} que coincidem em $\intcc{-1}1$
        coincidem em toda parte: \cref{cor:b1:poly:nroots}.)
  \item Deduza as \omterm{def:b1:structures:law}{leis de composição} e de produto:
        \[
        T_m \circ T_n = T_{mn},
        \qquad
        2\,T_m T_n = T_{m+n} + T_{\abs{m-n}} .
        \]
  \item Recorde do \cref{exo:b1:poly:10} as raízes $x_k =
        \cos\frac{(2k+1)\pi}{2n}$ e os pontos de equioscilação
        $y_k = \cos\frac{k\pi}n$, com $T_n(y_k) = (-1)^k$. Escreva
        a fatoração completa de $T_n$ sobre $\R$ e
        justifique que os $y_k$ se intercalam: $y_n < x_{n-1} < y_{n-1}
        < \dots < x_0 < y_0$.
  \item Demonstre que $T_n(\cosh t) = \cosh(nt)$ para todo $t \in \R$
        (mesma indução, usando a
        \cref{prop:b1:functions:hyprules}) e deduza, para $x
        \geq 1$, a forma fechada
        \[
        T_n(x) = \frac{\bigl(x + \sqrt{x^2 - 1}\bigr)^n +
        \bigl(x - \sqrt{x^2 - 1}\bigr)^n}{2} ,
        \]
        de modo que $T_n(x) > 1$ para $x > 1$: fora de $\intcc{-1}1$ o
        \omterm{def:b1:poly:def}{polinômio} escapa imediatamente.
\end{enumerate}

\medskip
\noindent\textbf{Parte II --- Coeficientes, a família $U_n$ e uma
equação diferencial.}
\begin{enumerate}
  \item A partir da fórmula de De Moivre
        (\cref{cor:b1:complex:demoivre}), demonstre a expressão
        explícita
        \[
        T_n(x) = \sum_{0 \leq 2j \leq n} \binom{n}{2j}\,
        x^{\,n-2j}\,(x^2 - 1)^j ,
        \]
        e verifique-a para $n = 3$.
  \item Calcule $T_n(1)$, $T_n(-1)$ e $T_n(0)$ para todo $n$.
  \item Defina $U_n$ (\emph{de segunda espécie}) por $U_0 = 1$, $U_1 =
        2X$, $U_{n+1} = 2X\,U_n - U_{n-1}$. Demonstre que
        $U_n(\cos\theta) = \frac{\sin(n+1)\theta}{\sin\theta}$
        para $\theta \notin \pi\Z$, e que $T_n' = n\,U_{n-1}$
        para $n \geq 1$.
  \item Demonstre que $\abs{\sin n\theta} \leq n\,\abs{\sin\theta}$
        para todo $\theta$ (indução) e deduza a estimativa do tipo
        Markov
        \[
        \abs{T_n'(x)} \leq n^2
        \quad\text{em} \intcc{-1}1,
        \qquad\text{com} T_n'(\pm1) = (\pm1)^{n-1}\,n^2 .
        \]
  \item Mostre que $y = T_n$ satisfaz a equação diferencial
        \[
        (1 - x^2)\,y'' - x\,y' + n^2\,y = 0 ,
        \]
        derivando a identidade $\sin\theta\,
        T_n'(\cos\theta) = n\sin n\theta$ em relação a
        $\theta$; verifique diretamente para $T_2$.
\end{enumerate}

\medskip
\noindent\textbf{Parte III --- O teorema de extremalidade de Chebyshev.}
Seja $\widetilde T_n = 2^{1-n}\,T_n$ (\omterm{def:b1:poly:def}{mônico} pela questão 2) e
escreva $\norm{P}_\infty = \sup_{x \in \intcc{-1}1}\abs{P(x)}$.
\begin{enumerate}
  \item Justifique que $\norm{\widetilde T_n}_\infty = 2^{1-n}$, atingido
        com sinais alternados nos $n + 1$ pontos $y_n < \dots
        < y_0$.
  \item Suponha que algum $P$ \omterm{def:b1:poly:def}{mônico} de grau $n$ tivesse
        $\norm P_\infty < 2^{1-n}$, e ponha $D = \widetilde T_n -
        P$. Mostre que $\deg D \leq n - 1$ e que $D(y_k)$ tem o
        sinal estrito de $(-1)^k$ para cada $k = 0, \dots, n$.
  \item Deduza que $D$ tem pelo menos $n$ raízes reais distintas (uma
        em cada intervalo, pela propriedade do valor intermediário,
        usada aqui no nível do ensino médio e demonstrada no
        \cref{ch:b1:continuity}) e conclua o \emph{teorema de
        Chebyshev}: todo $P$ \omterm{def:b1:poly:def}{mônico} de grau $n$ satisfaz
        \[
        \norm{P}_\infty \geq 2^{1-n} .
        \]
  \item (Caso de igualdade, primeiro passo) Suponha agora $\norm P_\infty =
        2^{1-n}$ exatamente, com $P$ \omterm{def:b1:poly:def}{mônico} de grau $n$, e seja $D =
        \widetilde T_n - P$. Mostre que $(-1)^kD(y_k) \geq 0$ para
        todo $k$ e que, se $D(y_k) = 0$ num ponto \emph{interior}
        $y_k$ ($0 < k < n$), então $D'(y_k) = 0$ também.
        \emph{(Num $y_k$ interior, tanto $\widetilde T_n$ quanto
        $P$ atingem um extremo de valor absoluto $\norm{\cdot}
        _\infty$; uma função derivável tem derivada nula
        num extremo interior --- usado no nível do ensino médio,
        demonstrado no \cref{ch:b1:derivative}.)}
  \item (Caso de igualdade, conclusão) Conte as raízes de $D$ com
        \omterm{def:b1:poly:derivative}{multiplicidade} para mostrar que $D = 0$: o minimizante é
        \emph{único}, $P = \widetilde T_n$.
  \item Transporte para um segmento arbitrário $\intcc ab$: mostre que
        a norma do supremo mínima de um \omterm{def:b1:poly:def}{polinômio mônico} de grau $n$ em
        $\intcc ab$ é $2\bigl(\frac{b-a}4\bigr)^n$, atingida por um
        \omterm{def:b1:poly:def}{polinômio} de Chebyshev reescalonado. \emph{(Substitua $x =
        \frac{a+b}2 + \frac{b-a}2\,t$ e acompanhe o coeficiente
        líder.)}
\end{enumerate}

\medskip
\noindent\textbf{Parte IV --- Aplicações.}
\begin{enumerate}
  \item Trabalhe o caso $n = 3$ à mão: localize os extremos de
        $\widetilde T_3 = X^3 - \frac34X$ em $\intcc{-1}1$,
        verifique a equioscilação em quatro pontos com valor
        $\frac14$ e conclua que nenhuma cúbica \omterm{def:b1:poly:def}{mônica} faz melhor.
  \item (Nós ótimos de interpolação) Para $n + 1$ nós $x_0,
        \dots, x_n \in \intcc{-1}1$, o erro de interpolação é
        governado por $\omega(X) = \prod_i (X - x_i)$ (como o
        \cref{ch:b1:taylor} quantificará). Demonstre que a escolha
        que minimiza $\norm\omega_\infty$ é o \omterm{def:b1:logic:sets}{conjunto} das $n + 1$
        raízes de $T_{n+1}$, com $\norm\omega_\infty = 2^{-n}$:
        os nós de Chebyshev são os lugares certos para interpolar.
  \item Usando $T_5$, demonstre que $c = \cos 36^\circ$ satisfaz
        $16c^5 - 20c^3 + 5c + 1 = 0$, fatore esse \omterm{def:b1:poly:def}{polinômio} como
        $(x + 1)(4x^2 - 2x - 1)^2$ e conclua que
        \[
        \cos 36^\circ = \frac{1 + \sqrt5}4 .
        \]
        Verifique a coerência com $\cos 72^\circ =
        \frac{\sqrt5 - 1}4$ do \cref{exo:b1:complex:8}.
  \item Estime $T_{10}(1.1)$ com a forma fechada da questão 6
        (dois algarismos significativos bastam) e interprete: um
        \omterm{def:b1:poly:def}{polinômio} limitado por $1$ em $\intcc{-1}1$ já pode
        exceder $40$ em $x = 1.1$. (Que $T_n$ cresça \emph{o mais
        rápido} entre tais \omterm{def:b1:poly:def}{polinômios} é outra propriedade
        extremal da família, além deste problema.)
  \item Demonstre a \omterm{def:b1:arith:congruence}{congruência}: para todo \omterm{def:b1:arith:prime}{primo} ímpar $p$, todos os
        coeficientes de $T_p - X^p$ são divisíveis por $p$.
        \emph{(Use a questão 7 e $p \mid \binom p{2j}$ para $0 <
        2j < p$, da demonstração do
        \cref{thm:b1:arith:fermat}.)} Verifique em $T_3$ e $T_5$.
\end{enumerate}

\medskip
\noindent\textbf{Parte V --- Síntese.}
\begin{enumerate}
  \item Calcule explicitamente a quadrática \omterm{def:b1:poly:def}{mônica} de norma do supremo
        mínima em $\intcc01$ e o seu desvio. (Questão 17
        com $n = 2$.)
  \item Onde exatamente o problema usou: (i) a rigidez dos
        \omterm{def:b1:poly:def}{polinômios} (\cref{cor:b1:poly:nroots}); (ii) a
        trigonometria do \cref{ch:b1:complex} e do
        \cref{ch:b1:functions}; (iii) a aritmética dos \omterm{def:b1:counting:objects}{coeficientes
        binomiais} do \cref{ch:b1:arith}? Uma frase para cada.
  \item Síntese, num parágrafo curto: o teorema diz que
        o \omterm{def:b1:poly:def}{polinômio mônico} mais achatado é aquele que
        \emph{equioscila}, e a demonstração converte a otimalidade
        numa contagem de raízes. Comente esse mecanismo, o
        papel da substituição $x = \cos\theta$ como ponte
        entre álgebra e trigonometria, e nomeie os dois
        lugares em que o problema precisou de fatos de análise (TVI,
        extremo interior) que capítulos posteriores demonstram.
\end{enumerate}
\end{problem}
